Bridge failure, a context that holds the endpoints of a reasoning chain but not the passage linking them, is strongly associated with the answer-quality loss on bridge questions for every generator we tested (F1 drops by 48--57 points between complete and failed contexts, also within the same question). A connectivity-aware selector that uses no LLM call recovers part of that loss: on HotpotQA it gives about 10 more points of gold-complete retrieval and 4.0--6.5 more F1 points than Hybrid RAG, and on 2WikiMultiHopQA it leads all LLM-free baselines by 8--19 points of gold-complete retrieval. The study also delimits the claim. The spectral form of the connectivity reward is not necessary, the benefit requires a graph that links bridge passages, it does not transfer to MuSiQue's deeper chains, and iterative LLM retrieval remains clearly stronger when its cost is acceptable. The iterative retriever, which finds the bridge passage with LLM-written queries, shows that candidate generation remains the larger lever; a connectivity-aware selector cannot recover a bridge that the relevance score and graph never surface. Better bridge discovery without per-query LLM calls, learned or LLM-assisted selection over a connected candidate pool, and graph construction for corpora without explicit links are the natural next steps.
